\subsection*{同一历史里有几种计数}

重新从初始 10000、空业务记录出发。客户端发送 R1，服务 A 取得 S，事务形成 P1、O1、D1 并提交。响应丢失，客户端发送 R2；B 尝试插入 S 得到零行，读取 P1 并返回原结果。若 R2 响应也丢失，R3 再做同样恢复。这条历史中有三次到达服务的 HTTP 尝试，三次服务处理入口调用，一次进入实际扣款分支，一次成功业务提交和一份已提交扣款效果。

投递次数在这里指请求消息实际到达所观察接收点的次数。接收点必须写清：到网关三次，可能只有两次到购买进程；客户端调用三次，消息也可能有一次未到任何服务。执行次数同样需要指定对象。整个处理函数调用三次，数据库插入尝试三次，余额更新只执行一次。若把这些都简写为执行一次，观察日志时很容易误判协议已经失败。

提交次数也应指定事务种类。R2 可以在一个只读或未产生业务写入的事务中结束，数据库仍可能完成一次事务控制流程。这里统计购买业务提交，指使 S 的终态及相应业务记录共同成立的那次提交。查询事务结束、访问日志入库、重放计数更新都不增加 P1 的购买次数。生效次数则以选定业务效果计数，本例按已提交扣款流水及与之对应的账户变化判断。

\begin{figure}[htbp]
\centering
\begin{tikzpicture}[x=1cm,y=1cm]
\node[booknode,text width=3cm] (a) at (0,0) {R1：执行余额更新\\事务随后回滚};
\node[booknode,text width=3cm] (b) at (5,0) {R2：更新并提交\\形成 P1 / O1 / D1};
\draw[bookarrow] (a)--(b);
\node[booknote] at (2.5,1) {两次尝试仍使用原 S};
\node[booknote,text width=3.5cm] at (0,-2.0) {SQL 可执行两次\\回滚更新不形成扣款效果};
\node[bookresult,text width=4cm] at (5,-2.0) {成功：终态 1，扣款效果 1\\拒绝：终态 1，扣款效果 0};
\node[booknode,text width=9cm] at (2.5,-3.8) {另一合法 S2 可再形成一次扣款；各身份的计数分别计算};
\end{tikzpicture}
\caption{请求到达、语句执行、事务提交与业务生效使用不同的计数口径。}
\label{fig:ch05-counts}
\end{figure}


因此可以准确地说：在本章假设内，同一完整身份至多形成一份已提交终态；成功终态对应一份购买、订单和扣款，重复请求只恢复原结果。它不要求服务入口只运行一次，也不要求每条 SQL 只发送一次。若首次为余额不足，购买终态提交一次，成功扣款生效次数为零。把每一终态都数成一次扣款，会漏掉这个必要分支。

若需要用一个计数概括这项上限，可以把固定 S 对应的已提交终态数量记为 \(N_T(S)\)，把其成功扣款效果数量记为 \(N_D(S)\)。在本章单次全额扣款模型中，要求 \(N_T(S)\leq 1\)，并且成功终态对应 \(N_D(S)=1\)，拒绝终态对应 \(N_D(S)=0\)。不存在已提交终态时，按共同事务关系也应没有该 S 的已提交扣款。三次请求恢复 P1 的历史中，入口处理为 3 次，\(N_T(S)=1\)、\(N_D(S)=1\)；余额不足历史中则分别为 1 和 0。这些数值只针对固定身份，不把 Alice 的另一项合法 S2 合进来。

\subsection*{进入过扣款分支也可能没有已提交扣款}

另取同一初始状态。A 取得 S，执行余额更新，自己的事务内读到 7000；随后写订单出错，数据库回滚。B 继续原 S，插入成功，再执行条件更新并提交 P2、O2、D2。余额更新语句一共实际运行两次，事务内部也两度得到 7000，最后却只有 D2 与已提交余额 7000。A 的内部修改已经撤销，没有成为可恢复的 D1 资金事实。

日志应保留 A 曾尝试更新，便于查明订单错误，但财务核对不能把两条更新日志相加成支出 6000。相反，只统计最后余额也不够：同一 S 的错误实现若两次各扣 3000，会留下 4000；另一项合法 S2 也各扣 3000，正确实现同样得到 4000。应同时检查完整身份、终态、订单与流水关联，以及实际提交的金额。

现在让 A 发送 COMMIT 后失联，B 在唯一检查等待。若 A 最后提交，B 返回 P1，更新语句仍只有 A 的一次；若 A 回滚，B 才可能实际更新。这项判断由数据库事务结局决定。服务提前开始多少次恢复，并不能增加取得同一 S 的已提交资格，前提仍是所有路径都经过相同非空组合唯一约束和共同事务。

\subsection*{恰好一次需要同时说明什么}

至少一次投递通常表示传送方在没有取得足够确认时继续尝试，使消息可能多次到达；要保证最终到达，还需要网络恢复、接收方可用和持续重试。至多一次业务生效表示同一操作的效果数量有上限，允许一项操作一直没有成功。恰好一次业务生效还要求这项本来应该成功的操作最终确实成功一次。永久余额不足的请求已有明确拒绝终态，其成功扣款次数理应为零，因此不能把它纳入所有请求都必须扣一次的说法。

RIFL 所用的恰好一次远程调用语义，有稳定调用身份、自动重试、可靠完成记录、重试查找及客户端资格等前提。论文讨论的是在这些机制和故障范围中如何把重复通信接到同一完成结果\cite{rifl}。本书当前将保证限定为单一权威数据库、单付款路径、记录仍保留和许可仍有效，并且已提交事实不被恢复流程丢弃。外部支付或异步副本切换需要各自的执行与恢复规则。

更贴近本例的表达是分别说明三个条件。身份仲裁和共同事务限制重复终态及扣款；完成记录让同参恢复取得同一业务结果；在数据库与网络恢复、旧身份仍有效且调用方持续恢复等条件下，操作有机会完成并被客户端获知。前两项可以在暂时无法前进时继续成立，最后一项需要可用性和足够长的恢复机会。一次成功响应丢失以后，效果已经发生，获知结果仍可能需要很多次通信。

若要断言最终完成，还须规定：恢复者不会永久停止，数据库最终能处理这项事务，竞争与调度不会使它永久等不到执行，业务判断能够结束，并且允许恢复的期限足以容纳这些步骤。客户端最终获知还要求有一次完整响应能够到达并被持久保存。E1 采用固定截止时间，所以任意长断网都在保证之外；到期时仍未获知的操作转入业务事实核查，不能继续承诺旧提交入口必然返回终态。

还要区分完成一项操作与使购买成功。恢复者最终取得 S、执行判断并提交余额不足，也完成了该操作的终态处理；要进一步承诺成功扣款，还需有有效报价、足够余额和其他成功条件。如果另一项合法消费先取得账户余额，使本次条件不再成立，协议可以正确地产生拒绝结果。进展保证推动操作走到明确结论，具体结论仍由实际业务状态决定。

把 S 保存到客户端持久记录也属于这条链路。客户端在第一次发送后丢失 K1，只剩商品名和余额查询，就可能无法安全重建原操作。服务保存了完整结果，并不保证客户端总有足够定位信息。跨设备恢复可以通过账户内待确认操作列表提供原身份与权限，但列表必须来自服务已有的可靠记录；原请求尚未到达时，服务也无从列出客户端尚未发送成功的意图。

\subsection*{结果证据怎样接回完整执行历史}

到这里，完成记录对恢复所需的内容已经具体化：S 指定操作，原版本与输入决定如何比较，终态区分成功和业务拒绝，P1、O1、D1 与成交快照说明哪些业务事实成立，结果版本规定如何输出。当前余额与订单状态留在各自读取路径，本次 Date、R2 和追踪信息留在当前响应与诊断中。新键新意图可以合法买同一商品，旧键同参则回到原结论。

这组关系还要求逐个故障位置都兑现相同规则。在插入 S 以后退出，应当只留下回滚；在扣款以后、提交以前退出，也应回滚整组；提交成功但未返回时，应当能从另一连接找到完整终态；结果已经读回、响应再次丢失时，下一次恢复仍应得到同样的 P1。若某个出口留下已提交 \code{pending}、缺失结果或错误编号，协议描述就与执行历史不一致。

下一章将沿这些具体位置中断服务，再用原 S 恢复，核对账户、购买、订单、流水与返回结果；同时把每一步的允许状态写清，说明不重复生效、不同购买不被合并以及结果稳定分别依赖哪条规则。
\FloatBarrier
